% Complete Chapter 8 introduction, original PDF 107 / printed 97.
% block: further-directions-chapter-heading page:107 type:chapter
\clearpage\setcounter{section}{7}
\section{{Further directions and applications}}\label{chap:further-directions}

% block: further-directions-hodge-formulation page:107 type:ordinary-paragraph
{So far we have considered as our Hilbert complex the $L^2$ de Rham complex on a domain in $\mathbb R^n$, so that the resulting operator is the Hodge Laplacian. For $0\leq k\leq n$, the segment \eqref{eq:finite-element-forms-short-complex} gives rise to the weak formulation: given $f\in L^2\Lambda^k$, find $\sigma\in H\Lambda^{k-1}$, $u\in H\Lambda^k$, and $p\in\mathsf H^k$ such that}
\begin{equation}\label{eq:general-hodge-weak-formulation}
\begin{aligned}
\langle\sigma,\tau\rangle-\langle u,d\tau\rangle&=0,&\tau&\in H\Lambda^{k-1},\\
\langle d\sigma,v\rangle+\langle du,dv\rangle+\langle p,v\rangle&=\langle f,v\rangle,&v&\in H\Lambda^k,\\
\langle u,q\rangle&=0,&q&\in\mathsf H^k,
\end{aligned}
\end{equation}
{for which the corresponding strong formulation is given in \eqref{eq:finite-element-forms-hodge-strong}. We have shown stability, consistency, and convergence for the Galerkin solution to this problem using any of the four pairs of spaces in \eqref{eq:four-galerkin-form-pairs} and any polynomial degree $r\geq1$. The corresponding rates of convergence, for smooth solutions, are shown in Table~\ref{tab:galerkin-form-approximation-rates}.}

% block: further-directions-introduction page:107 type:ordinary-paragraph
{In this chapter, we quickly survey some other problems that can be addressed in the FEEC framework.}
